% !TEX root = main.tex
% English preamble: formatting, packages and macros.
% Chosen in main.tex with % !TEX root = main.tex
% English preamble: formatting, packages and macros.
% Chosen in main.tex with % !TEX root = main.tex
% English preamble: formatting, packages and macros.
% Chosen in main.tex with % !TEX root = main.tex
% English preamble: formatting, packages and macros.
% Chosen in main.tex with \input{preamble-en}. Its twin, preamble-es.tex, is
% identical except for the lines marked [language]: when you change one file,
% make the same change in the other.
% Document data (course, title, authors, emails) and the co-author notes switch
% live in main.tex. PACKAGES.md explains every package with examples.

% ------- Language and fonts ---------------------------------------
\usepackage[english]{babel}   % [language]
\usepackage[T1]{fontenc}       % correct hyphenation and copy-paste of accented words
\usepackage{lmodern}
\usepackage{microtype}         % subtle spacing improvements, fewer overfull lines

% ------- Margins and paragraphs -----------------------------------
\usepackage[top=25mm,
   bottom=20mm,
   left=20mm,
   right=20mm,
   headheight=15pt,            % room for the header (fancyhdr warns otherwise)
   marginparwidth=2cm]{geometry}
\usepackage{parskip}           % no indentation, blank space between paragraphs

% ------- Colors ---------------------------------------------------
% Load xcolor with its options before tikz, listings or hyperref,
% which load it without options and would cause an option clash.
\usepackage[x11names]{xcolor}

% ------- Mathematics ----------------------------------------------
\usepackage{amsmath}
\usepackage{mathtools}         % \coloneqq and other extensions to amsmath
\usepackage{amssymb}
\usepackage{mathrsfs}          % \mathscr

% physics and siunitx both define \qty. siunitx's version is kept:
% \qty{9.81}{m/s^2} writes a quantity with units. For physics' auto-sized
% brackets use \quantity(...), \pqty{...}, \bqty{...}, \abs{...}, \norm{...}.
\ExplSyntaxOn                  % hides the notice siunitx prints about this
\msg_redirect_name:nnn { siunitx } { physics-pkg } { none }
\ExplSyntaxOff
\usepackage{siunitx}           % \qty{9.81}{m/s^2}, \unit{kg}, aligned number columns
\usepackage{physics}           % \dv, \pdv, \abs, \norm, \ev, \mel, \comm, \vb, ...
\usepackage{braket}            % \bra, \ket, \braket{a|b}, \set{x | x>0}; overrides physics' \bra/\ket
\AtBeginDocument{\RenewCommandCopy\qty\SI}

% ------- Figures, tables and layout -------------------------------
\usepackage{graphicx}
\graphicspath{{figures/}}
\usepackage{float}             % [H] placement
\usepackage[labelfont=bf, labelsep=period]{caption}
\captionsetup[figure]{name=Fig.}
\usepackage{subcaption}
\usepackage{booktabs}          % \toprule, \midrule, \bottomrule
\usepackage{multicol}

% ------- Bibliography ---------------------------------------------
\usepackage{csquotes}
\usepackage[backend=biber, style=numeric-comp, sorting=none]{biblatex}
\addbibresource{refs.bib}

% ------- Header and footer ----------------------------------------
\usepackage{fancyhdr}
\pagestyle{fancy}
\fancyhf{}
\fancyhead[L]{\assignment}
\fancyhead[R]{\course}
\fancyfoot[R]{\thepage}
\renewcommand{\headrulewidth}{2pt}
\renewcommand{\footrulewidth}{1pt}

% First page: no header, author emails in the footer.
% Activated in main.tex with \thispagestyle{firstpage} right after \maketitle.
\fancypagestyle{firstpage}{%
   \fancyhf{}
   \fancyfoot[L]{\footnotesize\authoremails}
   \fancyfoot[R]{\thepage}
   \renewcommand{\headrulewidth}{0pt}
   \renewcommand{\footrulewidth}{1pt}%
}

% ------- Links and cross-references (keep last, in this order) ----
\usepackage[colorlinks, pdfusetitle]{hyperref}   % pdfusetitle: PDF metadata from \title and \author
\hypersetup{citecolor = Red1,
    linkcolor = blue,
    filecolor = magenta,
    urlcolor = red,
    anchorcolor = cyan}
\usepackage[capitalise]{cleveref}   % [language] \cref{fig:x} -> Fig. 1, \cref{eq:x} -> Eq. (1)

% ------- Author marks ---------------------------------------------
% \authormark{id}{symbol}: superscript after a name, links to that author's email.
% \authoremail{id}{symbol}{email}: link target with a clickable mailto: address.
\newcommand{\authormark}[2]{\texorpdfstring{\hyperlink{#1}{$^{#2}$}}{}}
\newcommand{\authoremail}[3]{\hypertarget{#1}{$^{#2}$\,\href{mailto:#3}{#3}}}

% ------- Commands -------------------------------------------------
\newcommand{\R}{\mathbb{R}}
\newcommand{\N}{\mathbb{N}}
\newcommand{\sub}{\subset}
. Its twin, preamble-es.tex, is
% identical except for the lines marked [language]: when you change one file,
% make the same change in the other.
% Document data (course, title, authors, emails) and the co-author notes switch
% live in main.tex. PACKAGES.md explains every package with examples.

% ------- Language and fonts ---------------------------------------
\usepackage[english]{babel}   % [language]
\usepackage[T1]{fontenc}       % correct hyphenation and copy-paste of accented words
\usepackage{lmodern}
\usepackage{microtype}         % subtle spacing improvements, fewer overfull lines

% ------- Margins and paragraphs -----------------------------------
\usepackage[top=25mm,
   bottom=20mm,
   left=20mm,
   right=20mm,
   headheight=15pt,            % room for the header (fancyhdr warns otherwise)
   marginparwidth=2cm]{geometry}
\usepackage{parskip}           % no indentation, blank space between paragraphs

% ------- Colors ---------------------------------------------------
% Load xcolor with its options before tikz, listings or hyperref,
% which load it without options and would cause an option clash.
\usepackage[x11names]{xcolor}

% ------- Mathematics ----------------------------------------------
\usepackage{amsmath}
\usepackage{mathtools}         % \coloneqq and other extensions to amsmath
\usepackage{amssymb}
\usepackage{mathrsfs}          % \mathscr

% physics and siunitx both define \qty. siunitx's version is kept:
% \qty{9.81}{m/s^2} writes a quantity with units. For physics' auto-sized
% brackets use \quantity(...), \pqty{...}, \bqty{...}, \abs{...}, \norm{...}.
\ExplSyntaxOn                  % hides the notice siunitx prints about this
\msg_redirect_name:nnn { siunitx } { physics-pkg } { none }
\ExplSyntaxOff
\usepackage{siunitx}           % \qty{9.81}{m/s^2}, \unit{kg}, aligned number columns
\usepackage{physics}           % \dv, \pdv, \abs, \norm, \ev, \mel, \comm, \vb, ...
\usepackage{braket}            % \bra, \ket, \braket{a|b}, \set{x | x>0}; overrides physics' \bra/\ket
\AtBeginDocument{\RenewCommandCopy\qty\SI}

% ------- Figures, tables and layout -------------------------------
\usepackage{graphicx}
\graphicspath{{figures/}}
\usepackage{float}             % [H] placement
\usepackage[labelfont=bf, labelsep=period]{caption}
\captionsetup[figure]{name=Fig.}
\usepackage{subcaption}
\usepackage{booktabs}          % \toprule, \midrule, \bottomrule
\usepackage{multicol}

% ------- Bibliography ---------------------------------------------
\usepackage{csquotes}
\usepackage[backend=biber, style=numeric-comp, sorting=none]{biblatex}
\addbibresource{refs.bib}

% ------- Header and footer ----------------------------------------
\usepackage{fancyhdr}
\pagestyle{fancy}
\fancyhf{}
\fancyhead[L]{\assignment}
\fancyhead[R]{\course}
\fancyfoot[R]{\thepage}
\renewcommand{\headrulewidth}{2pt}
\renewcommand{\footrulewidth}{1pt}

% First page: no header, author emails in the footer.
% Activated in main.tex with \thispagestyle{firstpage} right after \maketitle.
\fancypagestyle{firstpage}{%
   \fancyhf{}
   \fancyfoot[L]{\footnotesize\authoremails}
   \fancyfoot[R]{\thepage}
   \renewcommand{\headrulewidth}{0pt}
   \renewcommand{\footrulewidth}{1pt}%
}

% ------- Links and cross-references (keep last, in this order) ----
\usepackage[colorlinks, pdfusetitle]{hyperref}   % pdfusetitle: PDF metadata from \title and \author
\hypersetup{citecolor = Red1,
    linkcolor = blue,
    filecolor = magenta,
    urlcolor = red,
    anchorcolor = cyan}
\usepackage[capitalise]{cleveref}   % [language] \cref{fig:x} -> Fig. 1, \cref{eq:x} -> Eq. (1)

% ------- Author marks ---------------------------------------------
% \authormark{id}{symbol}: superscript after a name, links to that author's email.
% \authoremail{id}{symbol}{email}: link target with a clickable mailto: address.
\newcommand{\authormark}[2]{\texorpdfstring{\hyperlink{#1}{$^{#2}$}}{}}
\newcommand{\authoremail}[3]{\hypertarget{#1}{$^{#2}$\,\href{mailto:#3}{#3}}}

% ------- Commands -------------------------------------------------
\newcommand{\R}{\mathbb{R}}
\newcommand{\N}{\mathbb{N}}
\newcommand{\sub}{\subset}
. Its twin, preamble-es.tex, is
% identical except for the lines marked [language]: when you change one file,
% make the same change in the other.
% Document data (course, title, authors, emails) and the co-author notes switch
% live in main.tex. PACKAGES.md explains every package with examples.

% ------- Language and fonts ---------------------------------------
\usepackage[english]{babel}   % [language]
\usepackage[T1]{fontenc}       % correct hyphenation and copy-paste of accented words
\usepackage{lmodern}
\usepackage{microtype}         % subtle spacing improvements, fewer overfull lines

% ------- Margins and paragraphs -----------------------------------
\usepackage[top=25mm,
   bottom=20mm,
   left=20mm,
   right=20mm,
   headheight=15pt,            % room for the header (fancyhdr warns otherwise)
   marginparwidth=2cm]{geometry}
\usepackage{parskip}           % no indentation, blank space between paragraphs

% ------- Colors ---------------------------------------------------
% Load xcolor with its options before tikz, listings or hyperref,
% which load it without options and would cause an option clash.
\usepackage[x11names]{xcolor}

% ------- Mathematics ----------------------------------------------
\usepackage{amsmath}
\usepackage{mathtools}         % \coloneqq and other extensions to amsmath
\usepackage{amssymb}
\usepackage{mathrsfs}          % \mathscr

% physics and siunitx both define \qty. siunitx's version is kept:
% \qty{9.81}{m/s^2} writes a quantity with units. For physics' auto-sized
% brackets use \quantity(...), \pqty{...}, \bqty{...}, \abs{...}, \norm{...}.
\ExplSyntaxOn                  % hides the notice siunitx prints about this
\msg_redirect_name:nnn { siunitx } { physics-pkg } { none }
\ExplSyntaxOff
\usepackage{siunitx}           % \qty{9.81}{m/s^2}, \unit{kg}, aligned number columns
\usepackage{physics}           % \dv, \pdv, \abs, \norm, \ev, \mel, \comm, \vb, ...
\usepackage{braket}            % \bra, \ket, \braket{a|b}, \set{x | x>0}; overrides physics' \bra/\ket
\AtBeginDocument{\RenewCommandCopy\qty\SI}

% ------- Figures, tables and layout -------------------------------
\usepackage{graphicx}
\graphicspath{{figures/}}
\usepackage{float}             % [H] placement
\usepackage[labelfont=bf, labelsep=period]{caption}
\captionsetup[figure]{name=Fig.}
\usepackage{subcaption}
\usepackage{booktabs}          % \toprule, \midrule, \bottomrule
\usepackage{multicol}

% ------- Bibliography ---------------------------------------------
\usepackage{csquotes}
\usepackage[backend=biber, style=numeric-comp, sorting=none]{biblatex}
\addbibresource{refs.bib}

% ------- Header and footer ----------------------------------------
\usepackage{fancyhdr}
\pagestyle{fancy}
\fancyhf{}
\fancyhead[L]{\assignment}
\fancyhead[R]{\course}
\fancyfoot[R]{\thepage}
\renewcommand{\headrulewidth}{2pt}
\renewcommand{\footrulewidth}{1pt}

% First page: no header, author emails in the footer.
% Activated in main.tex with \thispagestyle{firstpage} right after \maketitle.
\fancypagestyle{firstpage}{%
   \fancyhf{}
   \fancyfoot[L]{\footnotesize\authoremails}
   \fancyfoot[R]{\thepage}
   \renewcommand{\headrulewidth}{0pt}
   \renewcommand{\footrulewidth}{1pt}%
}

% ------- Links and cross-references (keep last, in this order) ----
\usepackage[colorlinks, pdfusetitle]{hyperref}   % pdfusetitle: PDF metadata from \title and \author
\hypersetup{citecolor = Red1,
    linkcolor = blue,
    filecolor = magenta,
    urlcolor = red,
    anchorcolor = cyan}
\usepackage[capitalise]{cleveref}   % [language] \cref{fig:x} -> Fig. 1, \cref{eq:x} -> Eq. (1)

% ------- Author marks ---------------------------------------------
% \authormark{id}{symbol}: superscript after a name, links to that author's email.
% \authoremail{id}{symbol}{email}: link target with a clickable mailto: address.
\newcommand{\authormark}[2]{\texorpdfstring{\hyperlink{#1}{$^{#2}$}}{}}
\newcommand{\authoremail}[3]{\hypertarget{#1}{$^{#2}$\,\href{mailto:#3}{#3}}}

% ------- Commands -------------------------------------------------
\newcommand{\R}{\mathbb{R}}
\newcommand{\N}{\mathbb{N}}
\newcommand{\sub}{\subset}
. Its twin, preamble-es.tex, is
% identical except for the lines marked [language]: when you change one file,
% make the same change in the other.
% Document data (course, title, authors, emails) and the co-author notes switch
% live in main.tex. PACKAGES.md explains every package with examples.

% ------- Language and fonts ---------------------------------------
\usepackage[english]{babel}   % [language]
\usepackage[T1]{fontenc}       % correct hyphenation and copy-paste of accented words
\usepackage{lmodern}
\usepackage{microtype}         % subtle spacing improvements, fewer overfull lines

% ------- Margins and paragraphs -----------------------------------
\usepackage[top=25mm,
   bottom=20mm,
   left=20mm,
   right=20mm,
   headheight=15pt,            % room for the header (fancyhdr warns otherwise)
   marginparwidth=2cm]{geometry}
\usepackage{parskip}           % no indentation, blank space between paragraphs

% ------- Colors ---------------------------------------------------
% Load xcolor with its options before tikz, listings or hyperref,
% which load it without options and would cause an option clash.
\usepackage[x11names]{xcolor}

% ------- Mathematics ----------------------------------------------
\usepackage{amsmath}
\usepackage{mathtools}         % \coloneqq and other extensions to amsmath
\usepackage{amssymb}
\usepackage{mathrsfs}          % \mathscr

% physics and siunitx both define \qty. siunitx's version is kept:
% \qty{9.81}{m/s^2} writes a quantity with units. For physics' auto-sized
% brackets use \quantity(...), \pqty{...}, \bqty{...}, \abs{...}, \norm{...}.
\ExplSyntaxOn                  % hides the notice siunitx prints about this
\msg_redirect_name:nnn { siunitx } { physics-pkg } { none }
\ExplSyntaxOff
\usepackage{siunitx}           % \qty{9.81}{m/s^2}, \unit{kg}, aligned number columns
\usepackage{physics}           % \dv, \pdv, \abs, \norm, \ev, \mel, \comm, \vb, ...
\usepackage{braket}            % \bra, \ket, \braket{a|b}, \set{x | x>0}; overrides physics' \bra/\ket
\AtBeginDocument{\RenewCommandCopy\qty\SI}

% ------- Figures, tables and layout -------------------------------
\usepackage{graphicx}
\graphicspath{{figures/}}
\usepackage{float}             % [H] placement
\usepackage[labelfont=bf, labelsep=period]{caption}
\captionsetup[figure]{name=Fig.}
\usepackage{subcaption}
\usepackage{booktabs}          % \toprule, \midrule, \bottomrule
\usepackage{multicol}

% ------- Bibliography ---------------------------------------------
\usepackage{csquotes}
\usepackage[backend=biber, style=numeric-comp, sorting=none]{biblatex}
\addbibresource{refs.bib}

% ------- Header and footer ----------------------------------------
\usepackage{fancyhdr}
\pagestyle{fancy}
\fancyhf{}
\fancyhead[L]{\assignment}
\fancyhead[R]{\course}
\fancyfoot[R]{\thepage}
\renewcommand{\headrulewidth}{2pt}
\renewcommand{\footrulewidth}{1pt}

% First page: no header, author emails in the footer.
% Activated in main.tex with \thispagestyle{firstpage} right after \maketitle.
\fancypagestyle{firstpage}{%
   \fancyhf{}
   \fancyfoot[L]{\footnotesize\authoremails}
   \fancyfoot[R]{\thepage}
   \renewcommand{\headrulewidth}{0pt}
   \renewcommand{\footrulewidth}{1pt}%
}

% ------- Links and cross-references (keep last, in this order) ----
\usepackage[colorlinks, pdfusetitle]{hyperref}   % pdfusetitle: PDF metadata from \title and \author
\hypersetup{citecolor = Red1,
    linkcolor = blue,
    filecolor = magenta,
    urlcolor = red,
    anchorcolor = cyan}
\usepackage[capitalise]{cleveref}   % [language] \cref{fig:x} -> Fig. 1, \cref{eq:x} -> Eq. (1)

% ------- Author marks ---------------------------------------------
% \authormark{id}{symbol}: superscript after a name, links to that author's email.
% \authoremail{id}{symbol}{email}: link target with a clickable mailto: address.
\newcommand{\authormark}[2]{\texorpdfstring{\hyperlink{#1}{$^{#2}$}}{}}
\newcommand{\authoremail}[3]{\hypertarget{#1}{$^{#2}$\,\href{mailto:#3}{#3}}}

% ------- Commands -------------------------------------------------
\newcommand{\R}{\mathbb{R}}
\newcommand{\N}{\mathbb{N}}
\newcommand{\sub}{\subset}
